% ============================================================
%  FRFP Case Study (Clinical):
% IBM Watson for Oncology
% ============================================================

\section{FRFP Case Study (Clinical): IBM Watson for Oncology}
\label{app:frfp-watson-oncology}

This case study analyses IBM Watson for Oncology (WFO) through the FRFP lens: explicit semantics
in \(E\), tacit clinical endorsement in \(T\), run-level dynamics, boundary design, safe-horizon
constraints, and institutional operators. The aim is not to adjudicate clinical correctness
case-by-case, but to diagnose protocol structure: where tacit endorsement was asked to carry more
load than the boundary and governance mechanisms could support.

\subsection{Background}

After Watson’s \emph{Jeopardy!} success in 2011, IBM pursued healthcare as a major application
area and partnered with Memorial Sloan Kettering Cancer Center (MSKCC) to develop Watson for
Oncology (WFO) as a decision-support product: given a structured patient description, WFO would
return ranked treatment options with rationales, drawing on guideline representations, literature
abstractions, and MSK-derived training inputs.

Two clarifications matter for protocol analysis. First, multiple “Watson + oncology” efforts existed;
MD Anderson’s \emph{Oncology Expert Advisor} (OEA) was a separate Watson-based initiative that
reportedly consumed roughly \$62M and was terminated before use on real patients. Second, public
reporting (2017--2018) described cases in which WFO produced recommendations characterized by
internal reviewers and customers as unsafe or inconsistent with established standards; we treat
those reports as evidence of boundary and governance stress, not as the only basis for the FRFP
diagnosis.

IBM marketed WFO to hospitals globally, often positioning it as a way to deliver MSK-like expertise
in settings with fewer oncology specialists. In parallel, MD Anderson pursued OEA. Reported outcomes
of that effort (cost, delays, termination) are a useful institutional analogue because they reveal how
large budgets, vendor relationships, and audit triggers can substitute for (or displace) protocol-level
governance.

In 2017--2018, investigative reporting based on internal IBM documents and interviews described
multiple examples where outputs were flagged as unsafe or incorrect by reviewers and customers.
One commonly cited illustration involved a recommendation that conflicted with major safety warnings
for the relevant drug class in the stated clinical context. Hospitals also reported low local concordance
and “workflow mismatch,” and some partnerships were discontinued or scaled back.

FRFP use of these facts is narrow: they show (i) explicit semantics not robust to real-world coverage
gaps and heterogeneity; and (ii) boundary conditions under which endorsement could drift toward
deference (especially in low-specialist environments), making run-level harm plausible even if
per-episode failure were “rare.”

This trajectory makes WFO a natural FRFP case: there is a crisp split between explicit artefacts
(structured cases, guideline representations, ranked options, rationales) and tacit clinical expertise;
persistent boundary ambiguity between “advice” and endorsement; long-horizon reliance risks; and
institutional operators (vendor marketing, procurement, audit, committee oversight) that shaped
which runs were admissible.

\subsection{Explicit and Tacit Spaces in Watson for Oncology}

\subsubsection*{Explicit space \(E_{\mathrm{clinical}}\)}

In FRFP terms, Watson for Oncology’s explicit domain comprises:

\begin{itemize}
  \item structured patient representations (stage, biomarkers, comorbidities);
  \item codified guidelines (e.g.\ NCCN recommendations) and literature
        abstractions as ingested into Watson;
  \item curated training cases and MSK-derived examples, including substantial use of “synthetic”
        cases in public descriptions of the training approach (a key coverage and external-validity
        risk for \(E\));
   \item ranked treatment options with rationales and evidence links (where provided);
   \item any product-layer constraints (templates, allowed regimens, score thresholds) that shape outputs.
 \end{itemize}
\paragraph{Run object and update loops.}
For this case study the relevant object is not a single recommendation, but a family of runs:
repeated case entry \(\to\) ranked outputs \(\to\) clinician interaction \(\to\) institutional uptake, policy
updates, and eventual scale or rollback. Model and knowledge-base updates, guideline revisions, and
site-specific configuration changes are part of the run dynamics, not “outside the system.”

All of these are explicit artefacts manipulable by the AI system:
given a patient input, WFO transforms explicit encodings into explicit
recommendations via model internals and expert rules.

\subsubsection*{Tacit space \(T_{\mathrm{clinical}}\)}

Tacit space in this setting consists of:

\begin{itemize}
  \item individual oncologists' clinical expertise, including pattern
        recognition for atypical cases and nuanced trade-offs
        (toxicity vs.\ benefit, patient preferences, comorbidities);
  \item knowledge of local formularies, resource constraints, and institutional
        norms;
  \item shared, evolving community sense of what constitutes acceptable
        interpretation of evidence in specific cancer subtypes;
  \item ethical and legal sense of responsibility for patient outcomes.
\end{itemize}

Oncologists do not simply read guidelines; they interpret them.
They adjudicate between trials, reconcile contradictory evidence, and adapt
decisions to specific patient contexts in ways not fully codified in
explicit rules.

\subsection{Intended Boundary Design vs.\ De Facto Boundary}

\subsubsection*{Stated design: AI as decision support}

Publicly, IBM framed WFO as a clinical decision support tool:

\begin{itemize}
  \item the oncologist remains the decision-maker;
  \item WFO proposes options; the human reviews and chooses;
  \item the system encodes MSK expertise and guidelines but does not replace
        human judgement.
\end{itemize}

In FRFP terms, this corresponds to:
AI confined to explicit space \(E_{\mathrm{clinical}}\),
with a boundary at the \emph{clinician’s acceptance} (or rejection) of
a recommendation, mediated by their tacit state in \(T_{\mathrm{clinical}}\).

\subsubsection*{Marketing and perception: boundary drift}

However, several features of deployment and marketing blurred that boundary:

\begin{itemize}
  \item IBM and some partners promoted WFO as offering ``best possible
        treatment'' based on leading cancer centres, which implicitly
        encourages clinicians and administrators to treat its outputs as
        authoritative rather than tentative.
  \item At some hospitals (e.g.\ Jupiter Medical Center), executives appear
        to have acquired WFO partly for marketing value, positioning it as a
        differentiating capability, while clinicians later reported that the
        product was largely unusable.
  \item Early MD Anderson efforts focused more on ambitious promises and
        integration challenges than on rigorous, prospective validation
        in real patients; the system never reached routine clinical use
        before being shelved.
\end{itemize}

From an FRFP perspective, these patterns erode the explicit/tacit boundary:

\begin{itemize}
  \item explicit artefacts (ranked options) are framed as near-final answers;
  \item tacit endorsement risks becoming a rubber stamp on system outputs;
  \item institutional incentives push toward de facto automation, even where
        formal documentation still describes WFO as ``advisory''.
\end{itemize}

\subsection{Protocol-Level Failures in FRFP Terms}

We now describe the main failure modes of Watson for Oncology as failures of
FRFP protocol design, rather than isolated model bugs.

\subsubsection*{Z.1 Boundary Confusion and Over-Automation Pressure}

\paragraph{Problem.}
WFO blurred who was really ``deciding'':

\begin{itemize}
  \item The UI presented a small, ranked list of options with categorical
        labels, rather than a broad landscape of possibilities.
  \item Marketing emphasised Watson’s ``expertise'' and partnership with elite
        centres, encouraging deference.
  \item Internal documents show that some customers perceived it as
        insufficiently useful or safe, but this was discovered only after
        deployment and real-world testing.
\end{itemize}

In FRFP language:

\begin{itemize}
  \item the \emph{boundary} between explicit suggestion and tacit endorsement
        was not structurally protected;
  \item nothing in the protocol enforced that clinicians treat WFO outputs
        as provisional hypotheses requiring independent re-grounding;
  \item institutional messaging encouraged treating WFO’s explicit artefacts
        as if correctness lived in the system rather than in human tacit
        states.
\end{itemize}

\paragraph{Consequence.}
This boundary confusion made it plausible, in principle, that clinicians
under time pressure could over-rely on WFO for guidance, especially in less
specialised settings.
Even where clinicians ultimately rejected WFO in practice (as at several
sites), the FRFP design left no robust guardrail to prevent over-automation.

\subsubsection*{Z.2 Explicit-Space Pathologies: Training and Coverage}

\paragraph{Problem.}
Subsequent reporting and IBM internal documents revealed that:

\begin{itemize}
  \item WFO’s training relied heavily on small sets of synthetic or
        hypothetical cases labelled by one or a few MSK physicians, rather
        than large, diverse, real-patient cohorts;
  \item the case library lacked statistical design, undermining coverage and
        generalisability;
  \item recommendations sometimes deviated from widely adopted national
        guidelines (e.g.\ NCCN) or reflected ``unconventional'' interpretations
        of evidence;
  \item some recommendations were directly unsafe, e.g.\ suggesting a drug
        explicitly contraindicated for patients with active bleeding.
\end{itemize}

In FRFP terms, these are failures in explicit space \(E_{\mathrm{clinical}}\):

\begin{itemize}
  \item the mapping from patient encodings to recommended options was
        underspecified and poorly validated;
  \item the training set did not define a reliable basis for a semantics
        that clinicians could safely treat as approximating guideline-based
        reasoning;
  \item explicit representations (case abstractions and labelled options)
        did not faithfully capture the tacit standards of the oncology
        community.
\end{itemize}

\paragraph{Consequence.}
Even if boundaries had been perfectly designed and strictly respected, the
explicit artefacts were structurally fragile:
a human relying on them as ``MSK in a box'' would be misled about their coverage and limits. Within FRFP, this is an explicit-semantics failure: the system’s
machine-side representation does not reliably track the clinical distinctions it purports to encode,
so the artefacts handed to clinicians are not a dependable basis for endorsement. The downstream
risk is not only a wrong recommendation, but a workflow that can translate brittle artefacts into
real commitments i.e explicit layer fragility matters even with good boundary discipline.

\subsubsection*{Z.3 Safe-Horizon Violations and Degradation Risk}

\paragraph{Problem.}
Watson for Oncology was intended for repeated use across many patients and
cases:

\begin{itemize}
  \item clinicians would query WFO case after case, day after day;
  \item institutional partners hoped to build standardised WFO-informed
        workflows and even marketing narratives around it.
\end{itemize}

FRFP predicts that in such long-horizon workflows:

\begin{itemize}
  \item tacit grounding can degrade: clinicians may start to read fewer
        guidelines directly, rely more on WFO’s summaries, and accept its
        pattern of options as normal;
  \item hallucinations (unsafe or unsupported endorsements) become almost
        inevitable along sufficiently long runs if boundaries and re-grounding
        are not carefully designed.
\end{itemize}

Yet WFO’s design and deployments did not incorporate explicit safe-horizon
constraints:

\begin{itemize}
  \item there were no protocol-level limits on how many WFO-guided decisions
        could be made without deeper human re-grounding in primary evidence;
  \item there was no requirement for periodic re-validation of WFO’s
        recommendations against up-to-date guidelines for each cancer type;
  \item no explicit mechanisms enforced independent re-derivation of
        rationale by clinicians after a certain depth of WFO usage.
\end{itemize}

\paragraph{Consequence.}
In organisations that might have fully embraced WFO, FRFP would predict a
drift toward overreliance:
clinicians would likely start accepting WFO’s pattern of options as
``the way we treat here'', reducing their sensitivity to edge cases and
guideline changes.
In practice, the opposite happened at many sites: clinicians quickly lost
trust and utilisation collapsed.
From an FRFP perspective, this is almost a \emph{fortunate} failure mode:
trust broke before dangerous long-horizon degradation could fully set in.

\subsubsection*{Z.4 Institutional Operator Failures}

\paragraph{Problem.}
At the institutional level, WFO’s story involves:

\begin{itemize}
  \item high-level partnerships (MSK, MD Anderson, others) with ambitious
        goals and significant budgets;
  \item weak prospective clinical validation prior to marketing;
  \item internal IBM and hospital audits revealing cost overruns, lack of
        integration, and poor clinical utility;
  \item eventual sell-off of Watson Health assets, effectively ending the
        original oncology vision.
\end{itemize}

In FRFP’s institutional semantics, we can think of:

\begin{itemize}
  \item a population of tacit states \(T^{\mathrm{pop}}\) including clinicians,
        administrators, IBM staff, regulators, and patients;
  \item institutional operators \(G\) (e.g.\ hospital governance, vendor
        decision processes, regulatory expectations) that update these states
        in response to explicit artefacts (pilots, metrics, sales pitches,
        incident reports).
\end{itemize}

WFO’s failure suggests that:

\begin{itemize}
  \item institutional operators overweighted vendor demos and marketing
        narratives relative to rigorous FRFP-style admissibility checks
        (safe horizons, boundary clarity, explicit semantics);
  \item governance mechanisms were not designed to detect FRFP-structural
        problems (e.g.\ training on hypothetical cases, unsafe recommendations)
        early and decisively;
  \item in at least one major site (MD Anderson), an internal financial audit
        rather than an epistemic governance body drove termination.
\end{itemize}

\paragraph{Consequence.}
From an FRFP standpoint, this is an example of attempted \emph{explicit-only
governance}:
institutions relied on marketing materials, limited tests, and internal
dashboards rather than systematically designing boundaries, safe horizons,
and explicit/tacit roles.
The non-automation theorems in the main text suggest that such governance
cannot in general recover tacit correctness.

\subsection{Counterfactual: An FRFP-Aligned Watson for Oncology}

It is instructive to sketch what WFO might have looked like under a strong
FRFP design discipline.

\subsubsection*{Z.1* Explicit/Tacit Separation and UI}

\begin{itemize}
  \item WFO would be explicitly framed as an \emph{artefact generator}:
        it surfaces guideline passages, trial summaries, and similar cases,
        but does not present canonical ranked treatment lists as ``Watson’s
        recommendation''.
  \item The UI would highlight:
        \begin{itemize}
          \item multiple guideline-consistent options side by side;
          \item conflicting evidence or lack of coverage for the current case;
          \item explicit uncertainty markers (e.g.\ ``no close match in
                training set; exercise heightened caution'').
        \end{itemize}
  \item The system would avoid language suggesting that it ``knows'' the best
        treatment; instead, it would present itself as a \emph{search and
        summarisation engine} for explicit artefacts.
\end{itemize}

\subsubsection*{Z.2* Boundary Design and Documentation}

\begin{itemize}
  \item Boundary events would be clearly identified:
        signing a treatment plan, entering orders, or officially adopting a
        WFO-supported plan in the EHR.
  \item At each such boundary, the clinical interface would:
        \begin{itemize}
          \item show which WFO artefacts were consulted;
          \item show any divergences between WFO suggestions and base
                guidelines;
          \item require a short human justification for accepting any
                option that relies materially on WFO output.
        \end{itemize}
  \item Logs would treat boundary decisions as first-class events, enabling
        later audit of where WFO influenced care.
\end{itemize}

\subsubsection*{Z.3* Safe Horizons and Re-Grounding}

\begin{itemize}
  \item For each cancer type, safe horizons would cap how many WFO-assisted
        decisions could be made before:
        \begin{itemize}
          \item an independent human review of recent cases against up-to-date
                guidelines is required;
          \item WFO’s training and rules are re-validated on real cases from
                that institution.
        \end{itemize}
  \item For each clinician, usage statistics would highlight:
        \begin{itemize}
          \item how often WFO suggestions are accepted vs.\ overridden;
          \item whether there is a drift toward uncritical acceptance.
        \end{itemize}
  \item At an institutional level, committees would periodically review
        WFO’s impact on care patterns and outcomes, with authority to adjust
        horizons, restrict indications, or suspend use.
\end{itemize}

\subsubsection*{Z.4* Training and Explicit Semantics}

\begin{itemize}
  \item Training would rely on large, audited collections of real cases with
        diverse practice patterns and explicit probabilistic uncertainty,
        rather than small synthetic case sets.
  \item The system would maintain explicit versioned representations of
        guidelines and evidence, allowing clinicians to see exactly which
        sources underpin each suggestion.
  \item Changes in guidelines or evidence would create explicit version
        transitions in WFO’s internal knowledge, visible at the boundary.
\end{itemize}

\subsubsection*{Z.5* Institutional Governance}

\begin{itemize}
  \item A formal FRFP-style governance body (e.g.\ an AI clinical oversight
        committee) would:
        \begin{itemize}
          \item review WFO’s design before deployment;
          \item approve specific indications, safe horizons, and audit plans;
          \item require regular reporting on errors, overrides, and usage.
        \end{itemize}
  \item Marketing claims and external messaging would be constrained to
        match what the governance body has actually evaluated.
\end{itemize}

Such a counterfactual design cannot guarantee success, but it would align
WFO with FRFP’s structural constraints:
AI as explicit-only, correctness at human boundaries, finite safe horizons,
and explicit institutional accountability.

\subsection{Lessons for FRFP and Clinical AI}

The Watson for Oncology case supports five protocol-level claims (each already developed in the
opening Vol2 chapters and the Concept Lookup):
\begin{enumerate}
  \item \textbf{Explicit success does not imply tacit safety.} Demos and limited pilots do not certify
        coverage, external validity, or boundary robustness under real clinical heterogeneity.
  \item \textbf{Boundary rhetoric is not boundary design.} Labeling a tool “decision support” does not
        ensure that endorsement remains meaningfully tacit; boundary events must be engineered with
        required artefacts, logging, and authority.
  \item \textbf{Safe horizons are required for repeated clinical use.} Long-horizon deployment changes
        reliance patterns; horizons must cap ungrounded repetition and force periodic re-grounding.
  \item \textbf{Institutional non-automation is a safety condition.} Vendor metrics, dashboards, and
        marketing cannot substitute for governance operators empowered to constrain admissible runs.
  \item \textbf{Protocol vocabulary sharpens “overhype” diagnoses.} FRFP localizes failure in explicit
        semantics, boundary structure, horizon policy, and institutional operators, rather than in
        a generic narrative of execution quality.
\end{enumerate}

In this light, Watson for Oncology is not just an anecdote of an overhyped
AI.
It is an early, large-scale illustration of why a formal architecture like
FRFP is needed to reason about clinical AI systems whose capabilities and
institutional footprints will only grow.
